\section{Evaluation}
Im Rahmen dieses Projekts wurden viele funktionale und technische Anforderungen an das zu entwickelnde Visualisierungs-Framework definiert.
Bei genauer Betrachtung lässt sich sagen, dass die Anforderungen erfüllt wurden.

Eine dieser Anforderung war, das Framework so zu entwickeln,
dass es direkt in JupyterLab genutzt werden kann.
Die Anwendung lässt sich problemlos in Jupyter Notebooks einbinden, sodass diese Anforderung erfüllt worden ist.

Ein weiteres Kriterium war die interaktive Bedienbarkeit.
Durch die Verwendung von ipywidgets konnten benutzerfreundliche Bedienelemente wie Schieberegler und Dropdown-Menüs implementiert werden.
Diese ermöglichen eine intuitive Steuerung von Parametern, wie zur Anpassung von Rotationen und erfüllen somit die
Anforderung nach einer interaktiven Auseinandersetzung mit kinematischen Konzepten.

Die Kompatibilität mit SymPy wurde ebenfalls sichergestellt. Die Funktionen in Util.py können Matrixobjekte der Sympy Bibliothek entgegennehmen,
sofern die Symbole sich auf konkrete Werte zurückführen lassen.

Im inhaltlichen Fokus stand die Darstellung und Transformation von Rotationen im dreidimensionalen Raum.
Die Umsetzung verschiedener Darstellungsformen wie Eulerwinkel, Roll-Pitch-Yaw-Winkel, Rotationsmatrizen und Quaternionen wurde realisiert.
Körper können in dem Framework anhand jeder dieser Rotationdarstellungen transformiert werden.
Die entsprechenden Umrechnungen zwischen diesen Darstellungen wird durch Funktionen der Library util.py ermöglicht.

Auch die Simulation von radgetriebenen Robotern wurde integriert. Mit der Funktion move(robot, vel, steps)
lassen sich einfache Bewegungsabläufe in der Ebene darstellen. Das erweitert das Framework über reine Rotationen hinaus.

Die Verwendung von Poetry zur Verwaltung von Abhängigkeiten sorgt darüber hinaus für eine stabile und konsistente Entwicklungsumgebung,
was insbesondere für die langfristige Nutzung und Weiterentwicklung des Frameworks relevant ist.

Ein Aspekt, der im aktuellen Projekt\-status noch nicht voll\-stän\-dig um\-ge\-setzt wurde,
ist die Modellierung stationärer Roboter (z.B. Knickarmroboter) samt Vorwärts- und Rückwärtskinematik.
Allerdings wurde dies von Beginn an als zukünftige Erweiterung vorgesehen und fällt nicht in den unmittelbaren Umfang dieser Projektarbeit.

Es lässt sich feststellen, dass das Framework die definierten Anforderungen erfüllt.
Es bietet eine funktional umfangreiche, technisch stabile Plattform zur Unterstützung der Lehre im Bereich der robotischen Kinematik.
Die Grundlage für eine spätere Erweiterung um komplexere Robotermodelle wurde ebenfalls gelegt, womit auch die langfristigen Projektziele realistisch erscheinen.